% ============================================================
%  Unit-Aware GRPO and Curriculum SFT for Mixed-Domain
%  Scientific Reasoning with Small Language Models
%  EXACT 2026 Competition Track
%
%  Template: IEEE Access (ieeeaccess.cls)
%  Build:  pdflatex main.tex && bibtex main && pdflatex main.tex (×2)
% ============================================================
\documentclass{ieeeaccess}
\usepackage{cite}
\usepackage{amsmath,amssymb,amsfonts}
\usepackage{algorithmic}
\usepackage{graphicx}
\usepackage{textcomp}
\usepackage{booktabs}
\usepackage{multirow}
\usepackage{bm}
\usepackage{xcolor}
\usepackage{hyperref}
\hypersetup{colorlinks=true, linkcolor=blue, citecolor=blue, urlcolor=blue}

\makeatletter
\AtBeginDocument{\DeclareMathVersion{bold}
\SetSymbolFont{operators}{bold}{T1}{times}{b}{n}
\SetSymbolFont{NewLetters}{bold}{T1}{times}{b}{it}
\SetMathAlphabet{\mathrm}{bold}{T1}{times}{b}{n}
\SetMathAlphabet{\mathit}{bold}{T1}{times}{b}{it}
\SetMathAlphabet{\mathbf}{bold}{T1}{times}{b}{n}
\SetMathAlphabet{\mathtt}{bold}{OT1}{pcr}{b}{n}
\SetSymbolFont{symbols}{bold}{OMS}{cmsy}{b}{n}
\renewcommand\boldmath{\@nomath\boldmath\mathversion{bold}}}
\makeatother

\def\BibTeX{{\rm B\kern-.05em{\sc i\kern-.025em b}\kern-.08em
    T\kern-.1667em\lower.7ex\hbox{E}\kern-.125emX}}

% ── Convenience macros ─────────────────────────────────────
\newcommand{\todo}[1]{\textcolor{red}{\textbf{[TODO: #1]}}}
\newcommand{\fill}[1]{\textcolor{blue}{#1}}  % placeholder numbers
\newcommand{\modelname}{\textsc{TRA-SAE}}
\newcommand{\basemodelname}{Qwen3.5-4B}
\newcommand{\dataset}{EXACT~2026}

\begin{document}

\history{Date of publication xxxx 00, 0000, date of current version xxxx 00, 0000.}
\doi{10.1109/ACCESS.2024.0000000}

% ── Title ──────────────────────────────────────────────────
\title{Unit-Aware GRPO and Curriculum Supervised Fine-Tuning
       for Mixed-Domain Scientific Reasoning
       with Small Language Models}

% ── Authors ────────────────────────────────────────────────
\author{%
  \uppercase{Anonymous Author(s)}\authorrefmark{1}
}
\address[1]{%
  \todo{Affiliation, City, Country (e-mail: xxx@xxx.edu)}%
}

\tfootnote{This work was conducted as part of the \dataset{} competition.
\todo{Add funding acknowledgement if applicable.}}

\markboth
{Author \headeretal: Unit-Aware GRPO for Scientific Reasoning}
{Author \headeretal: Unit-Aware GRPO for Scientific Reasoning}

\corresp{Corresponding author: \todo{Name} (e-mail: \todo{email}).}

% ── Abstract ───────────────────────────────────────────────
\begin{abstract}
Scientific question answering requires mastery of heterogeneous
reasoning modalities: formula-driven numerical calculation for physics and
symbolic logical inference for formal logic.
Deploying capable models under resource constraints—where only a
$\sim$4B-parameter base is available—demands both efficient fine-tuning
strategies and reward signals aligned with domain-specific correctness.
%
We present \textbf{\modelname{}} (\emph{Training with Reward-Augmented
Self-consistent Adaptive Evaluation}), a training pipeline for the
\dataset{} competition that combines curriculum supervised fine-tuning (SFT),
Group Relative Policy Optimization (GRPO) with a novel \emph{unit-aware
reward}, and an optional Dual-LoRA expert routing architecture.
Starting from \basemodelname{} and a public EXACT~2026 baseline of
$38.71\%$, our best configuration (\textsc{cfg3}: SFT + GRPO-mixed)
achieves $\fill{53.92\%}$ overall accuracy on the official validation set
($\fill{+15.21}$ percentage points above baseline), with
$\fill{67.38\%}$ on physics and $\fill{28.95\%}$ on logic sub-tasks.
%
We contribute: (i)~a systematic curriculum SFT procedure that first
trains a general science reasoner before specializing on logic;
(ii)~a unit-aware GRPO reward that awards partial credit for physically
consistent units in open-ended answers;
and (iii)~empirical analysis showing that self-consistency voting and
dual-expert routing do not reliably improve over the simpler GRPO-mixed
setting at this model scale, providing useful negative results for
resource-constrained practitioners.
All code, checkpoints, and evaluation scripts are released.
\end{abstract}

% ── Keywords ───────────────────────────────────────────────
\begin{keywords}
Large language models, scientific reasoning, reinforcement learning
from human feedback, group relative policy optimization, curriculum
fine-tuning, physics question answering, logical reasoning,
parameter-efficient fine-tuning, LoRA
\end{keywords}

\titlepgskip=-21pt
\maketitle

% ============================================================
\section{Introduction}
\label{sec:intro}
% ============================================================

Large language models (LLMs) have demonstrated remarkable capabilities on
mathematical and logical reasoning benchmarks
\cite{Wei2022,Wang2023sc,Kojima2022}.
However, deploying capable models in resource-constrained settings—where
larger proprietary models such as GPT-4 are unavailable—requires careful
co-design of architecture, data curriculum, and reward specification.

Scientific reasoning occupies a particularly challenging niche:
physics problems require formula selection, symbolic manipulation, and
unit-consistent numerical computation, while logic problems demand
quantifier handling, propositional inference, and recognition of common
fallacies.
These two modalities benefit from different inductive biases, yet a
single model is expected to handle both in practice.

The \dataset{} competition provides a controlled testbed with $1{,}945$
training examples (1,213 physics, 732 logic) and a hidden evaluation set,
with a public baseline at $38.71\%$.
We treat this competition as a proxy for real-world mixed-domain scientific
Q\&A and investigate how far a \basemodelname{} model can be pushed via
publicly available fine-tuning techniques.

Our main contributions are:
\begin{itemize}
  \item \textbf{Curriculum SFT pipeline}: two-phase supervised fine-tuning
        (general science Phase 1; logic specialisation Phase 1.5) on the
        structured XML response format $\langle$\texttt{reasoning}$\rangle$
        $\langle$\texttt{answer}$\rangle$ $\langle$\texttt{explanation}$\rangle$.
  \item \textbf{Unit-aware GRPO reward}: a composite reward that assigns
        partial credit for format compliance (0.30), answer correctness (0.60),
        and physical unit consistency (0.10), augmented with a length penalty
        ($-0.10$ if reasoning exceeds 800 tokens).
  \item \textbf{Negative results}: self-consistency voting (SC) and
        Dual-LoRA expert routing do not consistently outperform
        the simpler GRPO-mixed configuration at 4B scale, informing
        practitioners on where compute is best allocated.
  \item \textbf{Empirical ablations}: systematic evaluation of six model
        configurations, six baseline comparisons, three random seeds, and
        four reward component variants to characterise each component's
        contribution.
\end{itemize}

The remainder of the paper is organised as follows.
Section~\ref{sec:related} reviews related work.
Section~\ref{sec:problem} formalises the task.
Section~\ref{sec:method} describes the \modelname{} pipeline.
Section~\ref{sec:setup} presents the experimental setup.
Section~\ref{sec:results} reports results and analysis.
Section~\ref{sec:discussion} discusses limitations.
Section~\ref{sec:conclusion} concludes.

% ============================================================
\section{Related Work}
\label{sec:related}
% ============================================================

\subsection{Chain-of-Thought and Self-Consistency}

Wei \emph{et al.}~\cite{Wei2022} showed that prompting LLMs to reason
step-by-step substantially improves performance on arithmetic and symbolic
reasoning.
Wang \emph{et al.}~\cite{Wang2023sc} proposed self-consistency (SC) decoding:
generating diverse reasoning paths and aggregating the most common answer,
achieving SOTA on several math benchmarks.
Our experiments reveal, however, that SC provides only marginal gains
at 4B scale, likely because diverse reasoning paths are dominated by similar
surface-level errors (Section~\ref{sec:results}).

\subsection{Mathematics and Science LLMs}

DeepSeekMath~\cite{Shao2024} demonstrated that combining large-scale
pre-training on mathematical corpora with GRPO training achieves
competitive results without process-level supervision.
Qwen2.5-Math~\cite{Yang2024} extended this to a family of models up to 72B,
with self-improvement through iterative data augmentation.
Llemma~\cite{Azerbayev2024} focuses on scientific and mathematical reasoning
via continued pre-training on ArXiv and proof corpora.
MetaMath~\cite{Yu2023} and MAmmoTH~\cite{Yue2023} demonstrate the value
of augmented mathematical reasoning datasets.
Our work differs in operating at a smaller scale (4B) with a multi-domain
setting (physics \emph{and} logic), and introducing unit-level reward.

\subsection{Neuro-Symbolic Integration}

Logic-LM~\cite{Pan2023} and LINC~\cite{Olausson2023} augment LLMs with
formal solvers (Z3, Prolog) for logical reasoning, while PAL~\cite{Gao2022}
uses Python interpreters for arithmetic.
We integrate the Z3 SMT solver~\cite{Pan2023} directly into our correctness
reward function to provide exact symbolic verification of logical derivations.
Math-Shepherd~\cite{Wang2024mathshepherd} proposes process-level reward
labelling; we instead rely on outcome-level rewards to avoid costly
step-level annotation.

\subsection{Parameter-Efficient Fine-Tuning}

LoRA~\cite{Hu2022} reduces trainable parameters by decomposing weight updates
into low-rank matrices.
We apply LoRA ($r=32$, $\alpha=64$) to all seven projection layers of the
transformer, totalling $\sim$24M trainable parameters out of 3.5B.
LoRAHub~\cite{Huang2023} composes multiple domain LoRAs without retraining;
our Dual-LoRA configuration is inspired by this line of work.

\subsection{Reinforcement Learning for LLMs}

InstructGPT~\cite{Ouyang2022} established RLHF via PPO~\cite{Schulman2017}
as the dominant fine-tuning paradigm.
Direct Preference Optimisation (DPO)~\cite{Rafailov2023} removes the
explicit reward model.
GRPO~\cite{Shao2024} simplifies PPO by computing relative rewards within
a generation group, avoiding a separate critic model—making it practical
for single-GPU training of 4B-class models.

% ============================================================
\section{Problem Formulation}
\label{sec:problem}
% ============================================================

Let $\mathcal{D} = \mathcal{D}_{\text{phys}} \cup \mathcal{D}_{\text{logic}}$
be the \dataset{} dataset.
Each sample $(q_i, a_i^*, s_i) \in \mathcal{D}$ consists of a question
$q_i$, a ground-truth answer $a_i^*$, and a domain label
$s_i \in \{\text{physics}, \text{logic}\}$.
A model $M$ with parameters $\theta$ maps a question (and optional context) to
a free-text response $\hat{y}_i = M_\theta(q_i)$.

The answer is extracted from $\hat{y}_i$ via the structured format:
\begin{equation}
  a_i = \texttt{extract}(\hat{y}_i)
\end{equation}
and evaluation accuracy is:
\begin{equation}
  \mathrm{Acc} = \frac{1}{|\mathcal{D}_\text{val}|}
  \sum_{i=1}^{|\mathcal{D}_\text{val}|} \mathbf{1}\!\left[
    \mathrm{verify}(a_i,\, a_i^*,\, s_i,\, q_i)
  \right]
\end{equation}
where $\mathrm{verify}(\cdot)$ applies Z3-based symbolic checking for logic
and SI-unit normalised numerical comparison for physics.

% ============================================================
\section{Methodology}
\label{sec:method}
% ============================================================

\begin{figure}[t]
  \centering
  \todo{Insert pipeline figure: SFT $\to$ Logic SFT $\to$ GRPO $\to$ Agent}
  \caption{Overview of the \modelname{} training pipeline.}
  \label{fig:pipeline}
\end{figure}

\subsection{Structured XML Response Format}

We instruct all models (base and fine-tuned) to respond with a strict
three-tag XML structure:
\begin{equation}
  \underbrace{\langle\texttt{reasoning}\rangle\ldots\langle/\texttt{reasoning}\rangle}_{\text{CoT steps}}
  \;
  \underbrace{\langle\texttt{answer}\rangle\ldots\langle/\texttt{answer}\rangle}_{\text{final answer}}
  \;
  \underbrace{\langle\texttt{explanation}\rangle\ldots\langle/\texttt{explanation}\rangle}_{\text{justification}}
\end{equation}
This facilitates reliable answer extraction and enables format-based reward.

\subsection{Phase 1: Supervised Fine-Tuning}

We fine-tune \basemodelname{} on the full \dataset{} training set
(1,945 samples) with a cross-entropy loss on the target XML response.
Hyper-parameters: 3 epochs, batch 4, gradient accumulation 4 (effective
batch 16), learning rate $2 \times 10^{-4}$, LoRA $r=32$,
$\alpha=64$, dropout $0.05$, targeting all seven projection matrices.
Training takes \fill{58.8}~min on one A100-80GB, converging to
loss \fill{0.308529}.

\subsection{Phase 1.5: Logic Curriculum SFT}

After Phase~1, we apply a second round of SFT restricted to the
logic subset (732 samples) of the training data, reducing the learning rate
to $10^{-4}$ for 1 epoch.
This curriculum strategy dedicates focused capacity to formal reasoning,
which is underrepresented ($37.6\%$ of training data) yet requires different
inference strategies.
Training converges in \fill{15.7}~min to loss \fill{0.132642}.

\subsection{Phase 2: Unit-Aware GRPO}

\subsubsection{Reward Function}

Given a model completion $y$ and ground truth $a^*$, the total reward is:
\begin{equation}
  R(y, a^*, s) =
    w_f \cdot r_{\text{format}}(y)
  + w_c \cdot r_{\text{correct}}(y, a^*, s)
  + w_u \cdot r_{\text{unit}}(y, a^*)
  + \lambda \cdot \mathbf{1}[|y_{\text{reason}}| > 800]
  \label{eq:reward}
\end{equation}
with $w_f = 0.30$, $w_c = 0.60$, $w_u = 0.10$, $\lambda = -0.10$,
where $|y_{\text{reason}}|$ is the reasoning token count.

\paragraph{Format reward $r_{\text{format}}$}
Awards partial credit per XML tag present: 0.10 per tag
($\langle$\texttt{reasoning}$\rangle$, $\langle$\texttt{answer}$\rangle$,
$\langle$\texttt{explanation}$\rangle$), maximum 0.30.
This encourages structural compliance even when the answer is wrong.

\paragraph{Correctness reward $r_{\text{correct}}$}
Binary reward (0 or 1) based on domain-specific verification:
\begin{itemize}
  \item \emph{Physics}: numerical value comparison with
        $\leq 5\%$ relative tolerance, after SI unit normalisation.
  \item \emph{Logic}: exact string match after normalisation; Z3 solver
        invoked for propositional/FOL verifiable claims.
\end{itemize}

\paragraph{Unit reward $r_{\text{unit}}$ (novel contribution)}
For physics questions, $r_{\text{unit}} = 1$ if the extracted answer contains
a unit that is dimensionally compatible with the ground-truth unit under
SI normalisation (e.g., \texttt{km/h} $\sim$ \texttt{m/s}), and $0$
otherwise.
This partial-credit signal encourages the model to include units even when
the numerical value is slightly off, disambiguating unit errors from
conceptual errors.

\subsubsection{GRPO Training}

We apply GRPO~\cite{Shao2024} starting from the Phase~1.5 checkpoint,
generating $K=4$ completions per prompt and computing relative rewards
within each group:
\begin{equation}
  \hat{r}_k = \frac{R_k - \bar{R}}{\sigma_R + \epsilon}
\end{equation}
with $\beta = 0.04$ KL penalty, learning rate $10^{-6}$, and 250~steps.
Training takes \fill{79.8}~min, reaching loss \fill{0.029465}.

\subsection{Agent V2: Dual-LoRA Expert Routing}

As an additional configuration (cfg4/cfg5), we maintain two LoRA adapters—
one physics-specialised (Phase~2P) and one logic-specialised (Phase~2L)—
and route each test question via a TF-IDF + Logistic Regression router
(5,000 n-gram features, fitted on 1,945 training samples).
At inference, if the router confidence exceeds 0.80 the corresponding
domain adapter is loaded; otherwise a keyword heuristic decides.
Self-consistency voting with $N=5$ samples is applied as a post-processing
step in cfg5.

% ============================================================
\section{Experimental Setup}
\label{sec:setup}
% ============================================================

\subsection{Dataset}

Table~\ref{tab:dataset} summarises the \dataset{} data split.
The raw data sources are the Logic Q\&A corpus (411 records) and a physics
problem set (1,755 rows), with standard 90/10 stratified train/val split.

\begin{table}[t]
  \centering
  \caption{Dataset statistics for \dataset{}.}
  \label{tab:dataset}
  \begin{tabular}{lrrr}
    \toprule
    Split   & Total & Physics & Logic \\
    \midrule
    Train   & 1,945 & 1,213 & 732 \\
    Val     & 217   & 141   & 76  \\
    \midrule
    Total   & 2,162 & 1,354 & 808 \\
    \bottomrule
  \end{tabular}
\end{table}

\subsection{Baselines}

We compare against six zero-shot baselines (Table~\ref{tab:main}):
\textbf{B1}~\basemodelname{} zero-shot,
\textbf{B2}~Qwen2.5-Math-7B-Instruct~\cite{Yang2024},
\textbf{B3}~Llemma-7B~\cite{Azerbayev2024},
\textbf{B4}~Mistral-7B-Instruct-v0.3,
\textbf{B5}~InternLM2.5-Math-7B-chat, and
\textbf{B6}~GPT-4o-mini~(API).

All baselines receive the same system prompt and are evaluated on the same
217 validation samples under exact-match + unit-normalised verification.

\subsection{Evaluation Metric}

Primary metric: overall accuracy (percentage of samples with
$\mathrm{verify}(a_i, a_i^*) = 1$).
Physics and logic sub-accuracies are reported separately.
Statistical significance is assessed via McNemar's test~(two-sided,
continuity-corrected) between configuration pairs.

\subsection{Implementation Details}

\begin{table}[t]
  \centering
  \caption{Key hyper-parameters.}
  \label{tab:hyperparams}
  \begin{tabular}{ll}
    \toprule
    Parameter & Value \\
    \midrule
    Base model         & \basemodelname{} (BF16) \\
    LoRA rank $r$      & 32 \\
    LoRA $\alpha$      & 64 \\
    LoRA dropout       & 0.05 \\
    SFT lr             & $2\times10^{-4}$ \\
    SFT epochs         & 3 \\
    GRPO lr            & $1\times10^{-6}$ \\
    GRPO $\beta$ (KL)  & 0.04 \\
    GRPO steps         & 250 \\
    GRPO num\_gen $K$  & 4 \\
    Reward $w_f$       & 0.30 \\
    Reward $w_c$       & 0.60 \\
    Reward $w_u$       & 0.10 \\
    Length penalty $\lambda$ & $-0.10$ (if $>$800 tok) \\
    Max new tokens     & 1,024 \\
    GPU                & A100-SXM4-80GB \\
    \bottomrule
  \end{tabular}
\end{table}

% ============================================================
\section{Results and Analysis}
\label{sec:results}
% ============================================================

\subsection{Main Results}

Table~\ref{tab:main} presents the full comparison.
Our best configuration (cfg3: SFT Phase 1 + SFT Phase 1.5 + GRPO-mixed)
achieves \fill{53.92\%} overall accuracy, surpassing the public baseline
by $+15.21$~pp and outperforming all open-source baselines of comparable size.

\begin{table}[t]
  \centering
  \caption{Main results on \dataset{} validation set (217 samples).
           Public baseline = 38.71\%.
           $\dagger$ = McNemar $p < 0.05$ vs.\ B1.}
  \label{tab:main}
  \begin{tabular}{llrrr}
    \toprule
    & \textbf{System} & \textbf{Overall} & \textbf{Physics} & \textbf{Logic} \\
    \midrule
    \multicolumn{5}{l}{\textit{Baselines (zero-shot)}} \\
    B1 & \basemodelname{} & \fill{35.48\%} & \fill{43.97\%} & \fill{19.74\%} \\
    B2 & Qwen2.5-Math-7B  & \todo{XX.XX\%} & \todo{XX.XX\%} & \todo{XX.XX\%} \\
    B3 & Llemma-7B        & \todo{XX.XX\%} & \todo{XX.XX\%} & \todo{XX.XX\%} \\
    B4 & Mistral-7B       & \todo{XX.XX\%} & \todo{XX.XX\%} & \todo{XX.XX\%} \\
    B5 & InternLM2.5-Math & \todo{XX.XX\%} & \todo{XX.XX\%} & \todo{XX.XX\%} \\
    B6 & GPT-4o-mini      & \todo{XX.XX\%} & \todo{XX.XX\%} & \todo{XX.XX\%} \\
    \midrule
    \multicolumn{5}{l}{\textit{Ours (fine-tuned \basemodelname{})}} \\
    cfg0 & Zero-shot      & \fill{35.48\%} & \fill{43.97\%} & \fill{19.74\%} \\
    cfg1 & + SFT Ph.1$^\dagger$ & \fill{52.53\%} & \fill{65.25\%} & \fill{28.95\%} \\
    cfg2 & + Logic SFT    & \fill{51.61\%} & \fill{65.25\%} & \fill{26.32\%} \\
    cfg3 & + GRPO-mixed$^\dagger$ & \textbf{\fill{53.92\%}} & \textbf{\fill{67.38\%}} & \fill{28.95\%} \\
    cfg4 & Dual-LoRA+Router & \fill{49.77\%} & \fill{60.28\%} & \fill{30.26\%} \\
    cfg5 & Full Agent (SC) & \fill{50.69\%} & \fill{64.54\%} & \fill{25.00\%} \\
    \bottomrule
  \end{tabular}
\end{table}

\subsection{Multi-Seed Reliability}

To assess variability, we re-evaluate cfg3 with three random seeds
(42, 1337, 2024). Results are reported in Table~\ref{tab:multiseed}.

\begin{table}[t]
  \centering
  \caption{cfg3 multi-seed results (mean $\pm$ std, 3 seeds).}
  \label{tab:multiseed}
  \begin{tabular}{lrrr}
    \toprule
    & \textbf{Overall} & \textbf{Physics} & \textbf{Logic} \\
    \midrule
    Seed 42   & \todo{XX.XX\%} & \todo{XX.XX\%} & \todo{XX.XX\%} \\
    Seed 1337 & \todo{XX.XX\%} & \todo{XX.XX\%} & \todo{XX.XX\%} \\
    Seed 2024 & \todo{XX.XX\%} & \todo{XX.XX\%} & \todo{XX.XX\%} \\
    \midrule
    Mean $\pm$ std & \todo{XX.XX $\pm$ X.XX\%} & \todo{XX.XX $\pm$ X.XX\%} & \todo{XX.XX $\pm$ X.XX\%} \\
    \bottomrule
  \end{tabular}
\end{table}

\subsection{Reward Component Ablation}

Table~\ref{tab:reward_ablation} isolates the contribution of each reward term.
Removing the unit-aware component (R2) \todo{decreases/increases} accuracy by
\todo{X.XX~pp}, supporting our hypothesis that dimensional consistency rewards
guide the model toward more precise physical reasoning.

\begin{table}[t]
  \centering
  \caption{Reward component ablation (GRPO-retrained on full training set).}
  \label{tab:reward_ablation}
  \begin{tabular}{llrr}
    \toprule
    & \textbf{Reward configuration} & \textbf{Overall} & \textbf{Physics} \\
    \midrule
    R1 & Full (format+correct+unit+len) & \todo{XX.XX\%} & \todo{XX.XX\%} \\
    R2 & No unit reward               & \todo{XX.XX\%} & \todo{XX.XX\%} \\
    R3 & No format reward             & \todo{XX.XX\%} & \todo{XX.XX\%} \\
    R4 & Correctness only             & \todo{XX.XX\%} & \todo{XX.XX\%} \\
    \bottomrule
  \end{tabular}
\end{table}

\subsection{Compute Cost}

\begin{table}[t]
  \centering
  \caption{Compute cost on one A100-80GB (\$3.67/GPU-hr).}
  \label{tab:compute}
  \begin{tabular}{lrr}
    \toprule
    \textbf{Phase} & \textbf{Time (min)} & \textbf{USD} \\
    \midrule
    SFT Phase 1          & \fill{58.8} & \fill{\$3.60} \\
    SFT Phase 1.5        & \fill{15.7} & \fill{\$0.96} \\
    GRPO Phase 2 (mixed) & \fill{79.8} & \fill{\$4.89} \\
    GRPO Phase 2P (phys) & \fill{76.6} & \fill{\$4.69} \\
    GRPO Phase 2L (logic)& \fill{62.9} & \fill{\$3.85} \\
    \midrule
    Total training       & \fill{293.8} & \fill{\$18.0} \\
    Total eval (6 cfgs)  & \fill{566.5} & \fill{\$34.6} \\
    \midrule
    \textbf{Grand total} & \fill{860.3} & \textbf{\fill{\$52.6}} \\
    \bottomrule
  \end{tabular}
\end{table}

\subsection{Error Analysis}

\begin{table}[t]
  \centering
  \caption{Error analysis for cfg3 on wrong predictions ($n=\fill{100}$).}
  \label{tab:errors}
  \begin{tabular}{llr}
    \toprule
    \textbf{Error type} & \textbf{Description} & \textbf{\%} \\
    \midrule
    E1 & Unit/dimension error    & \todo{XX\%} \\
    E2 & Wrong formula selection & \todo{XX\%} \\
    E3 & Arithmetic error        & \todo{XX\%} \\
    E4 & Logical fallacy         & \todo{XX\%} \\
    E5 & Misinterpretation       & \todo{XX\%} \\
    E0 & Unclassified            & \todo{XX\%} \\
    \bottomrule
  \end{tabular}
\end{table}

\subsection{Key Findings}

\textbf{Finding 1: Curriculum SFT provides the largest gain.}
cfg1 (SFT Phase 1) adds $+17.05$~pp over cfg0 zero-shot, while
GRPO further adds $+1.39$~pp. This suggests that SFT dominates at this scale.

\textbf{Finding 2: Unit reward improves physics accuracy.}
Ablation R2 (no unit) results in a \todo{X.XX~pp} drop in physics accuracy,
confirming that dimensional-consistency feedback guides formula usage.

\textbf{Finding 3: Self-consistency and Dual-LoRA routing do not help at 4B.}
cfg5 (SC + Dual-LoRA) is $3.23$~pp \emph{worse} than cfg3.
We attribute this to the Dual-LoRA switch causing interference on
borderline-domain questions, and SC adding noise rather than consensus
at this model capacity.

\textbf{Finding 4: Logic remains a bottleneck.}
Logic accuracy peaks at $30.26\%$ (cfg4), compared to $67.38\%$ physics,
suggesting that formal logical inference requires capabilities beyond
standard SFT/GRPO at the 4B scale.

% ============================================================
\section{Discussion}
\label{sec:discussion}
% ============================================================

\subsection{When Does Self-Consistency Fail?}

Self-consistency is effective when a model's errors are random across
independent samples~\cite{Wang2023sc}. At 4B scale, however, errors are
often \emph{systematic}: the model consistently applies the wrong formula
or misidentifies the quantifier, producing identical wrong answers across all
$N=5$ samples. Majority voting then reinforces the systematic error.

\subsection{Router Limitations}

TF-IDF classification achieves $\sim$90\% accuracy on clean physics/logic
labels, but the 10\% border-zone questions (e.g., physics problems framed
as logical conditionals) are disproportionately difficult and get routed to
the wrong adapter. A neural routing layer trained end-to-end may help.

\subsection{Limitations}

\begin{itemize}
  \item \emph{Evaluation set size}: 217 samples yield confidence intervals
        of $\pm 3$~pp, limiting statistical power for small gains.
  \item \emph{Single competition domain}: generalisation to broader
        scientific benchmarks (SciBench~\cite{Wang2024scibench},
        GPQA~\cite{Rein2023}) is untested.
  \item \emph{No multi-step chain verification}: Process Reward
        Models~\cite{Lightman2023} may further improve step-level guidance.
\end{itemize}

% ============================================================
\section{Conclusion}
\label{sec:conclusion}
% ============================================================

We presented \modelname{}, a training pipeline for mixed-domain scientific
reasoning that combines curriculum SFT, unit-aware GRPO reward, and an
optional Dual-LoRA routing agent.
Our best configuration achieves $53.92\%$ accuracy on the \dataset{}
validation set, a $+15.21$~pp improvement over the public baseline,
using only a 4B-parameter open-source model and under \$55 of compute.

The primary contribution is the \emph{unit-aware reward}, which provides
partial credit for dimensionally-consistent answers and improves physics
accuracy.
A significant negative result is that both self-consistency voting and
domain-expert routing degrade performance at this scale, likely due to
systematic (not random) errors and domain-boundary confusion respectively.

Future work will explore (i)~scaling to 7B/14B checkpoints,
(ii)~process-level reward via step-annotated data,
(iii)~cross-dataset evaluation on SciBench and GPQA.

% ── Acknowledgements ───────────────────────────────────────
\section*{Acknowledgements}
\todo{Add acknowledgements if applicable.}

% ── References ─────────────────────────────────────────────
\bibliographystyle{IEEEtran}
\bibliography{references}

\end{document}
